
En este capítulo se presentan los objetivos y el catálogo de requisitos del nuevo sistema informatico. Opcionalmente, se incluye el análisis de la brecha entre los requisitos planteados y la solución base seleccionada.

\section{Objetivos de Negocio}
La implantacion del sistema persigue una serie de objetivos de negocio vinculados a la mejora de la actividad diaria del distribuidor y a la relación operativa con las peluquerías clientes. A diferencia de los objetivos del sistema, centrados en las capacidades de la aplicación, estos objetivos expresan el impacto esperado sobre la forma de trabajar del negocio una vez la solución se encuentre en uso.

Los principales objetivos de negocio identificados son los siguientes:

\begin{itemize}
    \item Reducir la dependencia de procesos manuales y de herramientas dispersas en la gestión comercial, logistica y administrativa del distribuidor.
    \item Centralizar la información relevante del negocio para disminuir duplicidades, inconsistencias y perdida de contexto entre visitas, pedidos, repartos y cobros.
    \item Mejorar la eficiencia operativa del distribuidor, acortando el tiempo dedicado a tareas de coordinación, consulta y seguimiento.
    \item Facilitar el control de la actividad comercial diaria mediante una vision más estructurada de clientes, visitas, rutas y resultados.
    \item Aumentar la trazabilidad de las operaciones criticas del negocio, especialmente en acceso al sistema, gestión de usuarios, seguimiento de pedidos y registro de cobros.
    \item Mejorar la calidad del servicio ofrecido a las peluquerías clientes, permitiendoles consultar catálogo, pedidos e información de reparto desde un canal digital unificado.
    \item Favorecer la fidelizacion de los clientes profesionales mediante una relación más continua, más transparente y apoyada en herramientas digitales de valor operativo.
    \item Sentar una base organizativa y técnica que permita escalar la actividad comercial e incorporar nuevos usuarios, nuevos comerciales y nuevas funcionalidades sin depender de procedimientos informales.
\end{itemize}

En consecuencia, el valor de negocio esperado no se limita a digitalizar tareas ya existentes, sino a reordenar el flujo completo de información del distribuidor para hacerlo más controlable, más trazable y más preparado para futuras ampliaciones funcionales.

\section{Objetivos del Sistema}
El sistema desarrollado persigue la construcción de una plataforma web unificada que permita al distribuidor, a su red comercial y a los clientes profesionales operar sobre una misma base de información, evitando la fragmentación entre herramientas aisladas de catálogo, seguimiento comercial y gestión de pedidos.

De forma más concreta, los objetivos generales del sistema son los siguientes:

\begin{itemize}
    \item Centralizar en una única aplicación la gestión de usuarios, clientes, catálogo, visitas, rutas, pedidos y seguimiento de cobros.
    \item Garantizar un acceso diferenciado por roles, de manera que cada perfil utilice una interfaz y un conjunto de operaciones acordes a sus responsabilidades reales dentro del proceso comercial.
    \item Facilitar al distribuidor la administración operativa del sistema, incluyendo la gestión de usuarios, la auditoría de acciones relevantes, la consulta global de pedidos y la configuración transversal de ciertas políticas del sistema, como el rate limiting.
    \item Proporcionar al comercial una herramienta móvil y operativa para planificar su jornada, consultar la ficha de sus clientes, preparar pedidos, organizar repartos y registrar el resultado de sus visitas.
    \item Permitir al cliente profesional consultar el catálogo, revisar cartas de color, preparar pedidos, consultar su historial y conocer el estado logistico de sus entregas cuando exista un reparto real planificado.
    \item Mantener trazabilidad sobre los procesos criticos del sistema, especialmente en autenticación, gestión administrativa de usuarios, ciclo de vida del pedido y seguimiento de cobros.
    \item Integrar servicios externos de soporte cuando aporten valor claro al flujo funcional, como almacenamiento de imágenes, geocodificación, cálculo de rutas, generación de QR o soporte de instalación como aplicación web progresiva.
    \item Asegurar que la aplicación pueda evolucionar de forma incremental, manteniendo una base modular que permita ampliar funcionalidades futuras sin rehacer los módulos ya estabilizados.
\end{itemize}

\section{Metodologia de captura de requisitos}
\label{sec:RequirementsCaptureMethodology}
\input{./sections/development/requirements/RequirementsCaptureMethodology.tex}

\section{Division modular de requisitos}
\label{sec:ModularRequirements}
\section{División modular de requisitos}
\label{sec:ModularRequirements}

Con el objetivo de facilitar la comprensión del sistema y organizar las funcionalidades de forma coherente, los requisitos se agrupan en \textbf{módulos funcionales} que representan las principales áreas de responsabilidad de la aplicación. Cada módulo define una parte del dominio del problema, reúne funciones relacionadas y sirve como referencia para los \textbf{requisitos funcionales} y los \textbf{requisitos de información} que se detallan en los apartados posteriores.

La descomposición del sistema en módulos funcionales responde a los principios de diseño arquitectónico en ingeniería del \textit{software}, que promueven la organización del sistema en componentes con \textbf{responsabilidades bien definidas}, \textbf{alta cohesión} y \textbf{bajo acoplamiento}, con el objetivo de facilitar su comprensión, mantenimiento y evolución \citep[Ch.~6--7]{sommerville}.

\subsection{M1. Gestión de acceso, alta y autorización de usuarios}

Este módulo se encarga de la \textbf{identificación}, \textbf{alta}, \textbf{autenticación} y \textbf{autorización} de los usuarios del sistema. Su responsabilidad principal consiste en garantizar que cada persona accede a la aplicación con una cuenta válida y que las operaciones disponibles se ajustan al rol que tiene asignado.

Las funciones agrupadas en este módulo incluyen la creación de usuarios, la gestión de solicitudes de registro, el inicio y cierre de sesión, la recuperación o modificación de credenciales, la administración del perfil personal, el control del estado de las cuentas y la asignación de roles. También agrupa las comprobaciones de permisos necesarias para distinguir entre administradores, comerciales y clientes profesionales.

Además, el módulo incorpora funciones de \textbf{seguridad} y \textbf{trazabilidad} asociadas al acceso: registro de eventos de autenticación, auditoría de acciones administrativas sobre usuarios, bloqueo o desactivación de cuentas cuando corresponda y protección básica frente a abuso mediante políticas de \textit{rate limiting}. De esta forma, el resto de módulos parte de una base común de identidad, permisos y control de acceso.

\subsection{M2. Gestión comercial y clientes}

Este módulo se encarga de gestionar la relación entre el distribuidor, sus comerciales y los \textbf{clientes profesionales}. Su finalidad es centralizar las \textbf{fichas de las peluquerías cliente}, mantener actualizada su información de contacto, dirección, observaciones, horarios de atención, coordenadas geográficas y \textbf{asignación comercial}, de forma que el sistema conserve una visión ordenada de quién atiende a cada cliente y bajo qué condiciones operativas.

Sobre esa base, el módulo agrupa las funciones de \textbf{planificación comercial}: configuración de la jornada de cada comercial, definición de puntos de inicio y finalización, creación y seguimiento de \textbf{visitas}, ordenación de \textbf{rutas}, estimación de tiempos de desplazamiento y reorganización de la agenda cuando una visita no encaja en la ventana horaria prevista. Estas funciones conectan la gestión de clientes con los procesos logísticos, las notificaciones y los repartos asociados al módulo \texttt{M4}.

\subsection{M3. Catálogo, productos y coloración}

Este módulo se encarga de organizar y poner a disposición de los usuarios el \textbf{catálogo técnico y comercial} de productos profesionales. Su responsabilidad consiste en estructurar la oferta del distribuidor para que pueda consultarse de forma clara por administradores, comerciales y clientes profesionales.

Las funciones agrupadas en este módulo incluyen la gestión de \textbf{productos}, \textbf{categorías}, \textbf{líneas comerciales}, subcategorías, imágenes, descripciones, precios, estados de disponibilidad y recursos documentales asociados. También incluye la \textbf{consulta del catálogo} desde los perfiles comerciales y de cliente, de modo que la información de producto pueda utilizarse tanto en la venta como en la preparación de pedidos.

Dentro del mismo módulo se agrupa la gestión de \textbf{coloración}: cartas de color, referencias cromáticas, tonos, numeraciones, recursos visuales y relación opcional entre tonos y productos pedibles. Esta parte permite que la información técnica de color no quede separada del catálogo comercial, facilitando la consulta profesional y la conexión posterior con pedidos o fichas técnicas del salón.

\subsection{M4. Pedidos, entregas y cobros}

Este módulo se encarga de agrupar el \textbf{ciclo comercial de venta} desde la preparación de un pedido hasta su entrega y seguimiento económico básico. Su finalidad es que cliente profesional, comercial y administrador trabajen sobre un mismo flujo trazable, evitando que pedidos, repartos y cobros queden dispersos en canales externos.

\textbf{Las funciones de pedido} incluyen la creación de borradores, la selección de productos y referencias de color, la edición de líneas, la confirmación del pedido, la cancelación cuando proceda, la consulta del estado y la conservación de información relevante como importes, descuentos, observaciones, método de entrega y usuario responsable. El pedido actúa como unidad comercial principal y permite conocer qué se solicita, quién lo solicita y en qué situación se encuentra.

\textbf{Las funciones de entrega} agrupan la preparación de uno o varios repartos asociados a un pedido, la selección de cantidades servidas en cada entrega, el registro de bultos, la generación de etiquetas logísticas y la diferenciación entre reparto realizado por comercial y entrega delegada en agencia. Cuando el reparto depende de un comercial, el módulo se relaciona con la planificación de visitas de \texttt{M2}; cuando se delega en agencia, conserva la información necesaria para identificar el envío y aplicar el cargo correspondiente.

\textbf{Las funciones de cobro} permiten registrar pagos parciales o totales, indicar método de pago, fecha, notas y usuario responsable, así como calcular el importe pendiente y mantener un historial económico mínimo asociado al pedido. El módulo no pretende sustituir a un \textit{ERP} completo ni asumir la facturación avanzada, sino cubrir la \textbf{trazabilidad operativa} imprescindible para saber qué se ha pedido, qué se ha entregado y qué queda pendiente de cobrar.

\subsection{M5. Gestión técnica del salón y seguimiento de servicios}

Este módulo se encarga de proporcionar al cliente profesional herramientas para gestionar la \textbf{información técnica} de los servicios realizados en su propio salón. Su objetivo es trasladar al sistema parte del conocimiento que normalmente queda repartido entre notas internas, conversaciones, fotografías o recordatorios informales.

Las funciones agrupadas en este módulo incluyen la creación y mantenimiento de \textbf{fichas de clientes del salón}, el registro de \textbf{servicios realizados}, la definición de \textbf{fichas técnicas}, la selección de tonalidades, la anotación de fórmulas, observaciones, productos utilizados y resultados obtenidos. También contempla la organización del \textbf{historial de servicios} para que el profesional pueda consultar trabajos anteriores antes de repetir, ajustar o recomendar un tratamiento.

El módulo incorpora además funciones de apoyo profesional como plantillas reutilizables, documentación visual del resultado, sugerencias basadas en información previa y generación de borradores técnicos editables. Estas funciones permiten que la aplicación no se limite a vender producto, sino que también acompañe el uso técnico de esos productos en el contexto real de la peluquería.

\subsection{M6. Promociones, notificaciones y formaciones}

Este módulo se encarga de centralizar la \textbf{comunicación operativa y comercial} entre el distribuidor y sus clientes profesionales. Su finalidad es facilitar la difusión de información relevante sin depender exclusivamente de llamadas, mensajes aislados o comunicaciones no trazables.

\textbf{Las funciones de promociones} incluyen la creación, edición, publicación, segmentación y consulta de campañas comerciales, así como la asociación de descuentos o condiciones promocionales a productos o pedidos cuando corresponda. De este modo, el distribuidor puede comunicar acciones comerciales y el cliente profesional puede conocerlas desde la propia aplicación.

\textbf{Las funciones de formación} cubren la publicación de actividades formativas, la consulta de sesiones disponibles, la inscripción de usuarios y el seguimiento básico de participación. Junto a ello, el módulo agrupa las \textbf{notificaciones internas}, los recordatorios personales y los avisos relacionados con visitas, promociones, formaciones o eventos relevantes del sistema.

Cuando la comunicación requiere canales externos, el módulo contempla el uso de correo \textit{SMTP}, notificaciones \textit{Web Push} y soporte para el comportamiento propio de una \textit{PWA} \citep{nodemailersmtp2026,webpushapi2026,webdev-pwa}. Estas funciones no sustituyen la comunicación personal del distribuidor, pero proporcionan un canal organizado, persistente y alineado con el resto de procesos comerciales.

\subsection{M7. Administración transversal}

Este módulo se encarga de agrupar las capacidades de \textbf{gobierno} y \textbf{configuración} que afectan al sistema de forma transversal. Su responsabilidad no es representar un proceso comercial concreto, sino proporcionar al administrador herramientas para supervisar la operación, mantener parámetros comunes y controlar aspectos generales de funcionamiento.

\textbf{Las funciones administrativas} incluyen la consulta de auditoría, el seguimiento de acciones relevantes, la revisión de accesos, la exportación de información operativa cuando resulte necesario y la gestión de configuraciones globales. También agrupa ajustes que afectan a varios módulos, como canales de aviso, importes configurables, umbrales comerciales, estados auxiliares o criterios que condicionan el comportamiento general del sistema.

En el ámbito transversal, el módulo reúne la información necesaria para gestionar servicios auxiliares ya vinculados a la aplicación, conservar estados básicos de configuración y mantener parámetros comunes que puedan ser reutilizados por otros módulos.

El módulo también contempla funciones técnicas de soporte que no pertenecen a un único proceso de negocio, como políticas de protección frente a abuso, parámetros de ejecución, registros de operación y trazabilidad de comunicaciones. Estas capacidades permiten que la aplicación sea administrable y evolutiva sin mezclar la gestión diaria del negocio con detalles internos de implementación.

\subsection{M8. Funcionalidades adicionales}

Este módulo agrupa funcionalidades identificadas como posibles líneas de evolución del sistema. Su responsabilidad es recoger capacidades que amplían el valor de la aplicación, pero que no forman parte del núcleo operativo principal descrito por los módulos anteriores.

Las funciones previstas incluyen un \textbf{asistente inteligente de diagnóstico capilar} apoyado en el uso de un \textbf{capilógrafo} y en una presentación técnica de diagnóstico capilar cedida por la marca. La intención es trasladar ese proceso a la aplicación e integrarlo con inteligencia artificial para proponer tratamientos o productos a partir de mediciones, observaciones técnicas, historial del cliente del salón y validación profesional.

El módulo también contempla la \textbf{simulación visual de estilos o coloración mediante la cámara del móvil}, el reconocimiento de productos mediante cámara o imagen, la reserva \textit{online} para clientes finales vinculados a una peluquería profesional y un \textbf{asistente virtual tipo \textit{chatbot}} capaz de resolver dudas sobre el uso de la aplicación, el catálogo, los tratamientos capilares y los procesos comerciales.

Por último, \texttt{M8} incorpora como trabajo futuro la \textbf{automatización de servicios externos}, incluyendo posibles flujos con Factusol, \textit{n8n} u otros sistemas de tipo \textit{ERP} \citep{sdelsol-factusol,n8n-webhook,n8n-http-request}. Esta línea permitiría sincronizar clientes o artículos, generar albaranes o facturas a partir de pedidos y devolver a la aplicación estados administrativos, siempre que se definan las condiciones técnicas y operativas necesarias.

Este módulo permite separar las funcionalidades exploratorias del alcance principal del sistema, manteniendo documentadas sus necesidades funcionales e informacionales para que puedan abordarse en iteraciones posteriores sin alterar la estructura de los módulos centrales.


\section{Requisitos funcionales}
\label{sec:FunctionalRequirements}
A continuación, la Tabla~\ref{tab:functional-requirements-executive} presenta el catálogo de requisitos funcionales en formato ejecutivo. Los requisitos mantienen identificadores únicos para facilitar su trazabilidad, siguiendo las recomendaciones de identificación y seguimiento propuestas por ISO/IEC/IEEE 29148 \cite{ieee29148}. La versión extensa empleada durante el desarrollo se conserva en el repositorio como copia de apoyo.

{\scriptsize
\setlength{\tabcolsep}{3pt}
\renewcommand{\arraystretch}{1.13}
\newcommand{\RFModuleHeader}[1]{%
\hline\hline
\multicolumn{4}{|>{\RaggedRight\arraybackslash}p{0.94\linewidth}|}{\textbf{#1}}\\
\hline\hline
}
\begin{longtable}{|>{\RaggedRight\arraybackslash}p{0.11\linewidth}|>{\RaggedRight\arraybackslash}p{0.18\linewidth}|>{\RaggedRight\arraybackslash}p{0.18\linewidth}|>{\RaggedRight\arraybackslash}p{0.43\linewidth}|}
\caption{Catálogo ejecutivo de requisitos funcionales agrupado por módulos}
\label{tab:functional-requirements-executive}\\
\hline
\textbf{ID} & \textbf{Módulo} & \textbf{Actores} & \textbf{Requisito funcional sintetizado} \\
\hline
\endfirsthead
\hline
\textbf{ID} & \textbf{Módulo} & \textbf{Actores} & \textbf{Requisito funcional sintetizado} \\
\hline
\endhead
\RFModuleHeader{M1. Acceso, alta y seguridad de cuentas}
\texttt{RF-M1-01} & Acceso y alta & Cliente, administrador & Registrar solicitudes de alta y permitir la creación directa de cuentas, manteniendo revisión administrativa, aprobación, rechazo, activación, desactivación y trazabilidad de la gestión. \\
\hline
\texttt{RF-M1-02} & Acceso y alta & Todos los roles & Autenticar usuarios mediante credenciales seguras, permitir solo cuentas activas y autorizadas, limitar intentos abusivos y registrar accesos correctos y fallidos. \\
\hline
\texttt{RF-M1-03} & Acceso y alta & Administrador & Gestionar roles de administrador, comercial y cliente profesional, restringiendo navegación, pantallas y operaciones según el rol autenticado. \\
\hline
\texttt{RF-M1-04} & Acceso y alta & Administrador & Aplicar condiciones de autorización comercial sobre cuentas profesionales y conservar el estado operativo de cada usuario junto con su historial administrativo. \\
\hline
\texttt{RF-M1-05} & Acceso y alta & Todos los roles & Crear, mantener y cerrar sesiones de usuario, impidiendo la reutilización de sesiones expiradas o revocadas. \\
\hline
\texttt{RF-M1-06} & Acceso y alta & Todos los roles & Permitir consulta y edición del perfil propio, imagen, datos de contacto y cambio de credenciales conforme a las políticas de seguridad. \\
\hline
\RFModuleHeader{M2. Clientes profesionales, comerciales y rutas}
\texttt{RF-M2-01} & Clientes y comerciales & Administrador & Registrar clientes profesionales, actualizar su información operativa, asignarlos o reasignarlos a comerciales y conservar la trazabilidad de esas asignaciones. \\
\hline
\texttt{RF-M2-02} & Clientes y comerciales & Comercial, cliente & Consultar la ficha operativa del cliente con datos de contacto, localización, observaciones, asignación comercial y contexto de planificación. \\
\hline
\texttt{RF-M2-03} & Clientes y comerciales & Comercial, cliente & Planificar, registrar, completar, aplazar o cancelar visitas comerciales, notificando cambios relevantes y respetando la ventana horaria del cliente. \\
\hline
\texttt{RF-M2-04} & Clientes y comerciales & Comercial & Previsualizar y ordenar rutas comerciales diarias con geolocalización, estimación de tiempos y secuencia de visitas. \\
\hline
\texttt{RF-M2-05} & Clientes y comerciales & Comercial, cliente & Configurar la jornada diaria del comercial y mostrar al cliente una estimación de llegada cuando exista una visita de reparto real vinculada. \\
\hline
\RFModuleHeader{M3. Catálogo, productos y coloración}
\texttt{RF-M3-01} & Catálogo & Cliente, comercial & Consultar el catálogo activo de productos con información comercial, técnica e imagen asociada. \\
\hline
\texttt{RF-M3-02} & Catálogo & Cliente, comercial & Filtrar y buscar productos por categoría, línea, subcategoría, referencia, nombre y estado de disponibilidad. \\
\hline
\texttt{RF-M3-03} & Catálogo & Cliente, comercial & Visualizar cartas de color y referencias cromáticas asociadas a productos y líneas comerciales. \\
\hline
\texttt{RF-M3-04} & Catálogo & Administrador & Mantener productos, categorías, líneas, subcategorías, recursos de apoyo, cartas de color y referencias desde administración. \\
\hline
\RFModuleHeader{M4. Pedidos, reparto y cobro}
\texttt{RF-M4-01} & Pedidos y reparto & Cliente, comercial & Crear, editar, confirmar, cancelar y consultar pedidos, manteniendo líneas, productos, referencias de color, imágenes de apoyo, descuentos visibles y estados coherentes. \\
\hline
\texttt{RF-M4-02} & Pedidos y reparto & Comercial, cliente & Vincular pedidos confirmados a visitas de reparto, validar entregas mediante QR escaneado o introducción manual y sincronizar el estado logístico con la visita. \\
\hline
\texttt{RF-M4-03} & Pedidos y reparto & Comercial, administrador & Registrar el estado mínimo de cobro de pedidos entregados, incluyendo método, fecha, notas y usuario responsable. \\
\hline
\texttt{RF-M4-04} & Pedidos y reparto & Administrador, comercial, cliente & Consultar la trazabilidad cruzada de pedidos, repartos y cobros desde las vistas correspondientes a cada rol. \\
\hline
\RFModuleHeader{M5. Gestión técnica del salón}
\texttt{RF-M5-01} & Salón & Cliente profesional & Gestionar fichas de clientes finales del salón con datos de contacto, observaciones e historial asociado. \\
\hline
\texttt{RF-M5-02} & Salón & Cliente profesional & Registrar servicios realizados a clientes del salón con fecha, descripción, resultado y relación con la ficha correspondiente. \\
\hline
\texttt{RF-M5-03} & Salón & Cliente profesional & Registrar información técnica de tratamientos, fórmulas, tonalidades, tiempos, productos usados y observaciones profesionales. \\
\hline
\texttt{RF-M5-04} & Salón & Cliente profesional & Consultar el historial técnico de cada cliente del salón, filtrando servicios, tratamientos y resultados previos. \\
\hline
\texttt{RF-M5-05} & Salón & Cliente profesional & Registrar y consultar productos utilizados en servicios para facilitar seguimiento técnico y reposición. \\
\hline
\texttt{RF-M5-06} & Salón & Cliente profesional & Proponer sugerencias de productos a partir del historial de tratamientos y necesidades detectadas. \\
\hline
\texttt{RF-M5-07} & Salón & Cliente profesional & Gestionar plantillas reutilizables e imágenes del resultado final de los servicios. \\
\hline
\texttt{RF-M5-08} & Salón & Cliente profesional & Generar borradores de correo técnico a partir de la ficha del servicio para contactar con el cliente final. \\
\hline
\RFModuleHeader{M6. Comunicaciones, promociones y formación}
\texttt{RF-M6-01} & Comunicaciones & Administrador, cliente, comercial & Crear, editar, publicar, retirar y consultar promociones comerciales segmentadas por rol, vigencia, producto o cliente. \\
\hline
\texttt{RF-M6-02} & Comunicaciones & Administrador, cliente & Gestionar segmentos o categorías jerárquicas de clientes y aplicar promociones o descuentos por pertenencia heredada, resolviendo como máximo una promoción por línea de pedido. \\
\hline
\texttt{RF-M6-03} & Comunicaciones & Administrador, todos los roles & Emitir notificaciones internas, correos, recordatorios y avisos Web Push asociados a eventos relevantes del sistema. \\
\hline
\texttt{RF-M6-04} & Comunicaciones & Administrador, cliente, comercial & Publicar formaciones, mostrar su disponibilidad y permitir la inscripción o cancelación por parte de usuarios autorizados. \\
\hline
\RFModuleHeader{M7. Administración técnica e integraciones}
\texttt{RF-M7-01} & Administración técnica & Administrador & Agrupar las funciones administrativas reales: usuarios, clientes, catálogo, comunicaciones, pedidos, auditoría y ajustes de avisos automáticos. \\
\hline
\texttt{RF-M7-02} & Administración técnica & Administrador & Consultar información útil para la gestión diaria sin convertir los diagnósticos técnicos internos en pantallas de uso ordinario. \\
\hline
\texttt{RF-M7-03} & Administración técnica & Administrador & Mantener una base técnica de integraciones externas y servicios auxiliares que pueda ser reutilizada por scripts, validación e integraciones futuras. \\
\hline
\texttt{RF-M7-04} & Administración técnica & Administrador & Preparar la base de integración con sistemas de gestión comercial o ERP, especialmente para intercambios futuros con Factusol mediante automatización intermedia con n8n. \\
\hline
\texttt{RF-M7-05} & Administración técnica & Administrador & Generar propuestas de pedido a proveedor a partir de productos activos, criterios de reposición y datos persistentes de integración. \\
\hline
\texttt{RF-M7-06} & Administración técnica & Administrador & Consultar y exportar auditoría de accesos, acciones administrativas y eventos relevantes en formato trazable. \\
\hline
\texttt{RF-M7-07} & Administración técnica & Equipo técnico & Aplicar y diagnosticar internamente políticas globales de rate limiting sobre acceso y APIs sensibles, reservando su modificación a operación técnica. \\
\hline
\texttt{RF-M7-08} & Administración técnica & Administrador, cliente & Ofrecer soporte transversal de compatibilidad, PWA, navegador no soportado e instalación, manteniendo la comprobación detallada de servicios auxiliares como infraestructura interna. \\
\hline
\RFModuleHeader{\textcolor{red}{M8. Funcionalidades futuras}}
\textcolor{red}{\texttt{RF-M8-01}} & \textcolor{red}{Funcionalidades futuras} & \textcolor{red}{Cliente profesional} & \textcolor{red}{Proponer tratamientos mediante un asistente inteligente a partir de información técnica, historial y contexto del cliente del salón.} \\
\hline
\textcolor{red}{\texttt{RF-M8-02}} & \textcolor{red}{Funcionalidades futuras} & \textcolor{red}{Cliente profesional} & \textcolor{red}{Simular estilos o coloraciones sobre imágenes para apoyar decisiones técnicas futuras.} \\
\hline
\textcolor{red}{\texttt{RF-M8-03}} & \textcolor{red}{Funcionalidades futuras} & \textcolor{red}{Cliente profesional, comercial} & \textcolor{red}{Reconocer productos mediante cámara o imagen para consulta, pedido o soporte operativo.} \\
\hline
\textcolor{red}{\texttt{RF-M8-04}} & \textcolor{red}{Funcionalidades futuras} & \textcolor{red}{Cliente final, cliente profesional} & \textcolor{red}{Gestionar reservas online de clientes finales vinculadas a salones profesionales.} \\
\hline
\end{longtable}
}


\section{Requisitos no funcionales}
\label{sec:NonFunctionalRequirements}
Los requisitos no funcionales describen las propiedades de calidad y restricciones que debe cumplir el sistema más allá de las funcionalidades que proporciona. Para su clasificación se ha tomado como referencia el modelo de calidad definido en el estándar ISO/IEC 25010 \cite{iso25010}.

La \autoref{tab:non-functional-requirements-executive} resume los requisitos no funcionales en formato ejecutivo, agrupados por característica de calidad para facilitar su lectura y trazabilidad.

{\scriptsize
\setlength{\tabcolsep}{3pt}
\renewcommand{\arraystretch}{1.14}
\newcommand{\RNFCategoryHeader}[1]{%
\hline\hline
\multicolumn{3}{|>{\RaggedRight\arraybackslash}p{0.91\linewidth}|}{\textbf{#1}}\\
\hline\hline
}
\begin{longtable}{|>{\RaggedRight\arraybackslash}p{0.18\linewidth}|>{\RaggedRight\arraybackslash}p{0.15\linewidth}|>{\RaggedRight\arraybackslash}p{0.58\linewidth}|}
\caption{Catálogo ejecutivo de requisitos no funcionales agrupado por características de calidad}
\label{tab:non-functional-requirements-executive}\\
\hline
\textbf{Característica} & \textbf{ID} & \textbf{Requisito no funcional sintetizado} \\
\hline
\endfirsthead
\hline
\textbf{Característica} & \textbf{ID} & \textbf{Requisito no funcional sintetizado} \\
\hline
\endhead
\RNFCategoryHeader{RNF1. Usabilidad}
Usabilidad & \texttt{RNF-U-01} & Proporcionar una interfaz intuitiva que permita realizar las tareas principales con un número reducido de interacciones. \\
\hline
Usabilidad & \texttt{RNF-U-02} & Adaptar la interfaz a móviles, tabletas y ordenadores, manteniendo una experiencia de usuario consistente. \\
\hline
Usabilidad & \texttt{RNF-U-03} & Presentar la información de forma clara y estructurada para facilitar la comprensión de datos operativos, comerciales y administrativos. \\
\hline
Usabilidad & \texttt{RNF-U-04} & Mantener criterios básicos de accesibilidad: etiquetas en formularios, nombres accesibles, contraste suficiente, texto legible y estructura revisable mediante WAVE. \\
\hline
\RNFCategoryHeader{RNF2. Eficiencia de rendimiento}
Rendimiento & \texttt{RNF-P-01} & Responder en tiempos razonables bajo uso normal, priorizando login, catálogo, clientes y operaciones sobre pedidos. \\
\hline
Rendimiento & \texttt{RNF-P-02} & Optimizar la carga de recursos para permitir un funcionamiento fluido también en dispositivos móviles con capacidad limitada. \\
\hline
\RNFCategoryHeader{RNF3. Compatibilidad}
Compatibilidad & \texttt{RNF-C-01} & Ser compatible con navegadores web modernos e informar mediante una ruta específica cuando el navegador no cumpla el soporte mínimo. \\
\hline
Compatibilidad & \texttt{RNF-C-02} & Ser utilizable desde los sistemas operativos móviles más habituales, especialmente Android e iOS mediante navegador. \\
\hline
Compatibilidad & \texttt{RNF-C-03} & Integrarse con servicios externos necesarios: almacenamiento de imágenes, geocodificación, rutas, QR, SMTP transaccional y Web Push. \\
\hline
\RNFCategoryHeader{RNF4. Fiabilidad}
Fiabilidad & \texttt{RNF-F-01} & Mantener disponibilidad suficiente para los usuarios y minimizar interrupciones del servicio durante la operación ordinaria. \\
\hline
Fiabilidad & \texttt{RNF-F-02} & Garantizar la conservación correcta de los datos almacenados, evitando pérdidas o inconsistencias de información. \\
\hline
\RNFCategoryHeader{RNF5. Seguridad}
Seguridad & \texttt{RNF-S-01} & Restringir el acceso a funcionalidades según el rol autorizado de cada usuario autenticado. \\
\hline
Seguridad & \texttt{RNF-S-02} & Proteger los datos mediante autenticación, autorización, trazabilidad de accesos y validación de operaciones sensibles. \\
\hline
Seguridad & \texttt{RNF-S-03} & Incorporar rate limiting y validación de sesión para reducir accesos abusivos o repetitivos sobre login y API. \\
\hline
\RNFCategoryHeader{RNF6. Mantenibilidad}
Mantenibilidad & \texttt{RNF-M-01} & Organizar el sistema en módulos y capas técnicas que faciliten su mantenimiento, evolución y ampliación. \\
\hline
Mantenibilidad & \texttt{RNF-M-02} & Permitir actualizar o modificar funcionalidades sin afectar innecesariamente al funcionamiento global de la aplicación. \\
\hline
\RNFCategoryHeader{RNF7. Portabilidad}
Portabilidad & \texttt{RNF-PO-01} & Ejecutar la aplicación en distintos dispositivos mediante navegador web, sin requerir instalaciones nativas específicas. \\
\hline
Portabilidad & \texttt{RNF-PO-02} & Permitir instalación como aplicación web mediante manifiesto e iconos, y habilitar notificaciones Web Push cuando el navegador lo soporte. \\
\hline
Portabilidad & \texttt{RNF-PO-03} & Ofrecer soporte mínimo ante fallos puntuales de conectividad mediante \textit{service worker} y caché del \textit{app shell}. \\
\hline
\end{longtable}
}


\section{Requisitos de información}
\label{sec:InformationRequirements}
En esta sección se resumen los requisitos de información que el sistema debe almacenar, consultar y mantener. Para evitar duplicidad con el modelo conceptual de la \autoref{sec:conceptualmodel}, el detalle se presenta como diccionario compacto: entidad, atributos clave y relaciones principales.

{\scriptsize
\setlength{\tabcolsep}{3pt}
\renewcommand{\arraystretch}{1.12}
\newcommand{\RIModuleHeader}[1]{%
\hline\hline
\multicolumn{4}{|>{\RaggedRight\arraybackslash}p{0.93\linewidth}|}{\textbf{#1}}\\
\hline\hline
}
\begin{longtable}{|>{\RaggedRight\arraybackslash}p{0.08\linewidth}|>{\RaggedRight\arraybackslash}p{0.22\linewidth}|>{\RaggedRight\arraybackslash}p{0.40\linewidth}|>{\RaggedRight\arraybackslash}p{0.23\linewidth}|}
\caption{Diccionario ejecutivo de requisitos de información agrupado por módulos}
\label{tab:information-requirements-executive}\\
\hline
\textbf{Módulo} & \textbf{Entidad} & \textbf{Atributos clave} & \textbf{Relaciones principales} \\
\hline
\endfirsthead
\hline
\textbf{Módulo} & \textbf{Entidad} & \textbf{Atributos clave} & \textbf{Relaciones principales} \\
\hline
\endhead
\RIModuleHeader{M1. Acceso, usuarios y auditoría}
\texttt{RI-M1} & Usuario & Identificador, email, contraseña cifrada, nombre, teléfono, empresa, imagen, último acceso, fechas de creación y actualización. & Rol, estado de usuario, cliente profesional o comercial cuando aplique. \\
\hline
\texttt{RI-M1} & Rol & Código, nombre y descripción del perfil de acceso. & Usuarios asociados y reglas de autorización por interfaz. \\
\hline
\texttt{RI-M1} & Estado de usuario & Código y nombre del estado operativo de la cuenta. & Usuarios y solicitudes administrativas. \\
\hline
\texttt{RI-M1} & Solicitud de registro & Email, datos del solicitante, rol solicitado, origen, estado, fechas de solicitud y revisión, motivo de rechazo. & Estado de solicitud, origen, usuario revisor y usuario creado si se aprueba. \\
\hline
\texttt{RI-M1} & Catálogos de solicitud & Estados y orígenes controlados de las peticiones de alta. & Solicitudes de registro. \\
\hline
\texttt{RI-M1} & Registro de gestión de usuarios & Usuario afectado, administrador, acción, estado/rol anterior y nuevo, motivo, fecha. & Usuario objetivo, usuario administrador y tipo de acción. \\
\hline
\texttt{RI-M1} & Registro de acceso & Usuario, email intentado, tipo de evento, resultado, IP, agente, token de sesión, fecha. & Usuario, tipo de evento y tipo de resultado. \\
\hline
\RIModuleHeader{M2. Clientes profesionales, comerciales y rutas}
\texttt{RI-M2} & Cliente profesional & Usuario vinculado, comercial asignado, nombre, contacto, dirección, ciudad, provincia, geolocalización, ventana de visita y observaciones. & Usuario, comerciales, visitas, pedidos, segmentos y clientes del salón. \\
\hline
\texttt{RI-M2} & Comercial & Usuario vinculado, código, zona, datos operativos, horario y preferencias de ruta. & Usuario, clientes asignados, visitas y rutas. \\
\hline
\texttt{RI-M2} & Asignación comercial-cliente & Cliente, comercial, fecha de asignación, fecha de fin, motivo y usuario responsable. & Cliente profesional y comercial. \\
\hline
\texttt{RI-M2} & Visita comercial & Cliente, comercial, tipo, estado, fecha planificada, resultado, observaciones y contexto de replanificación. & Cliente, comercial, ruta, pedidos de reparto y notificaciones. \\
\hline
\texttt{RI-M2} & Catálogos de visita & Estados y tipos de visita comercial. & Visitas comerciales. \\
\hline
\texttt{RI-M2} & Ruta comercial & Comercial, fecha, estado, origen, destino, distancia, duración estimada y configuración de jornada. & Comercial, estado de ruta y visitas en ruta. \\
\hline
\texttt{RI-M2} & Visita en ruta & Ruta, visita, orden, ETA, distancia acumulada y duración estimada. & Ruta comercial y visita comercial. \\
\hline
\RIModuleHeader{M3. Catálogo, productos y coloración}
\texttt{RI-M3} & Producto & Nombre, referencia, descripción, categoría, línea, subcategoría, estado, imagen e información técnica. & Categoría, línea, subcategoría, estado, recursos, pedidos y productos usados en servicios. \\
\hline
\texttt{RI-M3} & Categoría de producto & Nombre, descripción, imagen y orden de visualización. & Líneas comerciales y productos. \\
\hline
\texttt{RI-M3} & Línea comercial & Categoría, nombre, descripción, imagen y orden. & Categoría, subcategorías, productos y cartas de color. \\
\hline
\texttt{RI-M3} & Subcategoría de producto & Línea, subcategoría padre, nombre, descripción, imagen y orden. & Línea, jerarquía de subcategorías y productos. \\
\hline
\texttt{RI-M3} & Recurso de apoyo & Título, descripción, tipo, URL, producto asociado y orden. & Producto y tipo de recurso. \\
\hline
\texttt{RI-M3} & Carta de color & Nombre, descripción, línea, imagen y orden. & Línea comercial y referencias de color. \\
\hline
\texttt{RI-M3} & Referencia de color & Carta, código, nombre, tono, equivalencia, imagen y orden. & Carta de color, líneas de pedido y fichas técnicas. \\
\hline
\texttt{RI-M3} & Catálogos de catálogo & Estados de producto y tipos de recurso. & Productos y recursos de apoyo. \\
\hline
\RIModuleHeader{M4. Pedidos, reparto y cobro}
\texttt{RI-M4} & Pedido & Cliente, comercial, estado, estado de cobro, importe, fechas, notas, QR de entrega y visita vinculada. & Cliente, comercial, estado, visita de reparto y líneas de pedido. \\
\hline
\texttt{RI-M4} & Línea de pedido & Pedido, producto, referencia de color, cantidad, precio y observaciones. & Pedido, producto y referencia cromática. \\
\hline
\texttt{RI-M4} & Estados de pedido y cobro & Código, nombre y descripción de estados operativos. & Pedidos. \\
\hline
\texttt{RI-M4} & Visita de reparto vinculada & Visita comercial que actúa como soporte logístico de entrega. & Pedido, cliente, comercial y ETA mostrada al cliente. \\
\hline
\RIModuleHeader{M5. Gestión técnica del salón}
\texttt{RI-M5} & Cliente del salón & Cliente profesional propietario, nombre, teléfono, email, observaciones y fechas. & Cliente profesional, servicios, sugerencias y reservas futuras. \\
\hline
\texttt{RI-M5} & Servicio realizado & Cliente del salón, fecha, tipo, descripción, resultado, observaciones e imágenes. & Cliente del salón, ficha técnica, productos usados e imágenes. \\
\hline
\texttt{RI-M5} & Ficha técnica del servicio & Diagnóstico, fórmula, tonalidad, tiempo de exposición, notas técnicas y resultado. & Servicio realizado y referencias de color. \\
\hline
\texttt{RI-M5} & Producto utilizado en servicio & Servicio, producto, referencia cromática, cantidad, uso y observaciones. & Servicio, producto y referencia de color. \\
\hline
\texttt{RI-M5} & Plantilla de servicio & Cliente profesional, nombre, descripción y configuración reutilizable. & Cliente profesional y productos de plantilla. \\
\hline
\texttt{RI-M5} & Imagen de resultado & Servicio, URL, descripción, orden y fecha de subida. & Servicio realizado. \\
\hline
\texttt{RI-M5} & Sugerencia de producto & Cliente del salón, producto, motivo, prioridad y fecha. & Cliente del salón y producto. \\
\hline
\texttt{RI-M5} & Correo técnico generado & Servicio, destinatario, asunto, contenido y fecha de generación. & Servicio y ficha técnica. \\
\hline
\RIModuleHeader{M6. Comunicaciones, promociones y formación}
\texttt{RI-M6} & Promoción & Título, descripción, tipo, beneficio, vigencia, estado, segmento y condiciones. & Productos, segmentos, clientes y pedidos con descuento. \\
\hline
\texttt{RI-M6} & Segmento de cliente & Código, nombre, criterio, descripción y fechas. & Clientes profesionales, promociones y asignaciones. \\
\hline
\texttt{RI-M6} & Asignación de segmento & Cliente, segmento, origen, fecha de asignación y usuario responsable. & Cliente profesional y segmento. \\
\hline
\texttt{RI-M6} & Producto promocionado & Promoción, producto y condición de aplicación. & Promoción y producto. \\
\hline
\texttt{RI-M6} & Notificación interna & Destinatario, título, mensaje, tipo, canal, estado de lectura y fechas. & Usuario, cliente, visita, promoción o formación según origen. \\
\hline
\texttt{RI-M6} & Suscripción push & Usuario, endpoint, claves, estado, agente, último uso y revocación. & Usuario y notificaciones Web Push. \\
\hline
\texttt{RI-M6} & Recordatorio & Usuario, título, mensaje, fecha programada, estado y origen. & Usuario y evento funcional relacionado. \\
\hline
\texttt{RI-M6} & Formación e inscripción & Datos de formación, fecha, modalidad, aforo, estado, usuario inscrito y estado de inscripción. & Formación, usuario e integración externa si aplica. \\
\hline
\RIModuleHeader{M7. Administración técnica e integraciones}
\texttt{RI-M7} & Configuración del sistema & Clave, valor, tipo, ámbito, descripción y fecha de actualización. & Usuario administrador y servicios transversales. \\
\hline
\texttt{RI-M7} & Política de rate limiting & Ámbito, límite, ventana temporal, estado y descripción. & Sobrescritura técnica y endpoints protegidos. \\
\hline
\texttt{RI-M7} & Integración externa & Nombre, tipo, estado, configuración, último chequeo y observaciones. & Operaciones de integración, comunicaciones y propuestas a proveedor. \\
\hline
\texttt{RI-M7} & Operación de integración & Integración, tipo de operación, tipo de dato, estado, payload, resultado y fechas. & Integración externa y datos intercambiados. \\
\hline
\texttt{RI-M7} & Propuesta de pedido a proveedor & Proveedor, estado, criterio de generación, importe estimado, fecha y observaciones. & Integración externa y líneas de propuesta. \\
\hline
\texttt{RI-M7} & Línea de propuesta a proveedor & Propuesta, producto, cantidad, motivo y datos de reposición. & Propuesta y producto. \\
\hline
\RIModuleHeader{\textcolor{red}{M8. Funcionalidades inteligentes previstas}}
\textcolor{red}{\texttt{RI-M8}} & \textcolor{red}{Consulta de recomendación} & \textcolor{red}{Usuario, cliente del salón, tipo y estado del cabello, tratamientos previos y observaciones.} & \textcolor{red}{Usuario, cliente del salón y recomendaciones generadas.} \\
\hline
\textcolor{red}{\texttt{RI-M8}} & \textcolor{red}{Recomendación de tratamiento} & \textcolor{red}{Consulta, producto sugerido, descripción, prioridad y motivo.} & \textcolor{red}{Consulta de recomendación y producto.} \\
\hline
\textcolor{red}{\texttt{RI-M8}} & \textcolor{red}{Simulación de coloración} & \textcolor{red}{Usuario, imagen origen, cambio solicitado, resultado y fecha.} & \textcolor{red}{Usuario y posible referencia de producto.} \\
\hline
\textcolor{red}{\texttt{RI-M8}} & \textcolor{red}{Reconocimiento de producto} & \textcolor{red}{Usuario, imagen, código detectado, producto, propósito y fecha.} & \textcolor{red}{Usuario, producto y propósito de reconocimiento.} \\
\hline
\textcolor{red}{\texttt{RI-M8}} & \textcolor{red}{Reserva online} & \textcolor{red}{Cliente profesional, cliente del salón, contacto, servicio, fecha y estado.} & \textcolor{red}{Cliente profesional, cliente del salón y estado de reserva.} \\
\hline
\end{longtable}
}


\section{Reglas de negocio}
Las reglas de negocio describen restricciones operativas que el sistema debe respetar para reflejar de manera coherente el funcionamiento esperado por el distribuidor y por su red comercial. En la implementación actual, las principales reglas de negocio identificadas son las siguientes:

\begin{itemize}
    \item Solo los usuarios autenticados, activos y autorizados pueden acceder a las funcionalidades protegidas del sistema.
    \item El acceso a rutas, vistas y operaciones queda condicionado por el rol del usuario, distinguiendo al menos entre administrador, comercial y cliente profesional.
    \item Las solicitudes de registro de nuevos clientes no suponen acceso inmediato: deben quedar pendientes de revisión administrativa hasta su aprobacion o rechazo.
    \item Un cliente profesional puede estar asignado a uno o varios comerciales según la estructura operativa definida por el distribuidor, pero un comercial solo puede operar directamente sobre clientes que tenga asignados.
    \item Las visitas comerciales deben quedar asociadas a un cliente existente y a un comercial responsable, y su estado debe evolucionar de manera coherente entre planificada, completada, aplazada o cancelada.
    \item La hora exacta de una visita o entrega no se trata como un dato fijo introducido manualmente, sino como una estimación dependiente de la planificación de ruta y del contexto operativo del día.
    \item El cliente solo debe visualizar una ETA cuando exista una visita de reparto real, valida y vinculada a pedidos coherentes para la jornada.
    \item Los pedidos pueden iniciarse como borrador y confirmarse posteriormente, pero no deben saltar arbitrariamente entre estados incompatibles.
    \item Un mismo cliente no puede acumular más de dos pedidos abiertos pendientes de cobro al mismo tiempo.
    \item Solo pueden vincularse a una visita de reparto los pedidos confirmados del cliente correcto; no deben poder asociarse pedidos cancelados, ya entregados o pertenecientes a otro cliente.
    \item Una visita de reparto no debe poder completarse sin pedidos vinculados ni sin la validación QR correspondiente a cada pedido entregado.
    \item El cierre correcto de una visita de reparto debe marcar como entregados los pedidos asociados, manteniendo sincronizado el estado logistico del pedido con el de la visita.
    \item Si una visita de reparto se aplaza o se cancela, el sistema debe evitar estados huerfanos o inconsistentes entre visita, pedido y ETA mostrada al cliente.
    \item El registro de cobro solo aplica al alcance mínimo definido en esta iteración: distinguir entre pedidos pendientes de cobro y pedidos cobrados, con metodo, fecha, notas y usuario responsable.
    \item La gestión de cobros no incluye por ahora pagos parciales, facturación ni conciliación contable avanzada, por lo que el sistema no debe simular dichas capacidades.
    \item Las políticas de rate limiting forman parte de la operación del sistema y pueden configurarse globalmente desde administración para proteger el acceso y determinadas APIs frente a uso abusivo.
\end{itemize}

\section{Análisis GAP}
\textcolor{red}{TODO} Esta sección es opcional y es de aplicación cuando en la revisión de la técnica se hayan localizado soluciones informaticas que permitan servir de base para el nuevo sistema software. En esta sección se debe identificar y medir las diferencias entre lo que proporcionan dichas soluciones y los requisitos definidos para el proyecto. El resultado de este análisis permitira identificar cuales de estos requisitos ya estan solventados total o parcialmente por los distintos sistemas y poder elegir uno de ellos.
